\section{Các nghiên cứu liên quan}
Bối cảnh của các Hệ thống Phát hiện Xâm nhập Mạng (NIDS) đã có những bước phát triển vượt bậc trong những năm gần đây, nhờ vào sự sẵn có của các bộ dữ liệu toàn diện như CICIDS2017. Một lượng lớn các tài liệu nghiên cứu trong khoảng 3--5 năm qua đã khám phá các phương pháp tiếp cận dựa trên dữ liệu khác nhau; tuy nhiên, việc xử lý hiệu quả sự mất cân bằng giữa các lớp (multi-class imbalance) và đảm bảo hiệu quả triển khai tại biên vẫn là những thách thức quan trọng.

\subsection{Học máy truyền thống và Kỹ thuật đặc trưng}
Các thuật toán Học máy (ML) truyền thống, bao gồm Cây quyết định (Decision Trees - DT), Random Forest (RF) và k-Lân cận gần nhất (k-NN), đã được áp dụng rộng rãi cho NIDS nhờ sự đơn giản về mặt toán học \cite{breiman2001random, sharafaldin2018toward}. Các nghiên cứu gần đây đã chuyển trọng tâm sang việc tiêu chuẩn hóa không gian đặc trưng và tăng cường tính giải thích được cho mô hình. Sarhan và cộng sự \cite{sarhan2022towards} đã chứng minh rằng việc sử dụng một tập hợp đặc trưng thu gọn và được chuẩn hóa có thể duy trì độ chính xác phát hiện cao trong khi giảm đáng kể chi phí tính toán. Hơn nữa, việc cung cấp tính minh bạch cho mô hình đã trở nên tối quan trọng; Le và cộng sự \cite{le2022classification} cũng như Wang và cộng sự \cite{wang2022interpretable} đã sử dụng phương pháp SHAP (SHapley Additive exPlanations) để giải thích các quyết định của các mô hình ensemble dựa trên cây, đảm bảo rằng các mô hình phát hiện không chỉ chính xác mà còn đáng tin cậy. Tuy nhiên, các bộ phân loại đơn lẻ thường gặp khó khăn với lưu lượng mạng có chiều cao và thường nhạy cảm với các phân phối dữ liệu cụ thể \cite{gharaee2022deep}.

\subsection{Học sâu và các kiến trúc lai}
Để nắm bắt các mối quan hệ phi tuyến tính phức tạp, các nhà nghiên cứu ngày càng áp dụng nhiều kiến trúc Học sâu (Deep Learning - DL). Mạng thần kinh tích chập (CNN) và các kiến trúc tái phát như LSTM đã cho thấy khả năng vượt trội trong việc trích xuất đặc trưng tự động \cite{zhou2021deep}. Những tiến bộ gần đây nhấn mạnh vào các cơ chế chú ý lai (hybrid attention mechanisms) và sự hợp nhất đa giai đoạn. Hu và cộng sự \cite{hu2024improved} đã đề xuất một mạng dư (residual network) cải tiến với cơ chế chú ý lai, đạt độ chính xác 99,79\% trên bộ dữ liệu CICIDS2017. Tương tự, Ji và cộng sự \cite{ji2025hybrid} đã giới thiệu một khung hợp nhất đặc trưng hai giai đoạn song song (GRU-ResCBANet) để nắm bắt cả các phụ thuộc không gian và thời gian, báo cáo độ chính xác 99,83\% cho bài toán phân loại nhị phân. 

\subsection{Học Ensemble và các khung tính toán dữ liệu lớn}
Học Ensemble đã trở thành một chuẩn mực hiện đại trong lĩnh vực an ninh mạng. Để xử lý khối lượng lớn lưu lượng mạng hiện nay, các khung tính toán phân tán như Apache Spark \cite{zaharia2010spark} đã được tích hợp với các mô hình ensemble. Gupta và cộng sự \cite{gupta2022big} đã chứng minh rằng việc kết hợp các bộ phân loại học máy trong khung Spark cho phép phát hiện xâm nhập có khả năng mở rộng mà không làm giảm độ chính xác. Song song đó, Khan và cộng sự \cite{khan2024balanced} đã chứng minh rằng các mô hình ensemble dựa trên cây kết hợp với kỹ thuật cân bằng SMOTE-Tomek có thể đạt độ chính xác lên tới 99,96\% trên CICIDS2017. Giải quyết vấn đề mất cân bằng lớp ở mức độ sinh dữ liệu, Ding và cộng sự \cite{ding2024tmg} đã đề xuất mô hình TMG-GAN, tổng hợp hiệu quả các mẫu tấn công thuộc lớp thiểu số với điểm F1 đạt 99,72\%.

\subsection{Giải quyết khoảng trống nghiên cứu: Độ chính xác so với Hiệu quả}
Hạn chế căn bản trong các tài liệu đã được đánh giá là giả định ngầm về nguồn lực tính toán không bị giới hạn. Trong khi các kiến trúc như kiến trúc được đề xuất bởi Nguyen và cộng sự \cite{nguyen2024apelid} tập trung vào các mô hình ensemble song song thời gian thực, thì việc tích hợp các mô hình ensemble có độ chính xác cao với kỹ thuật nén mô hình phục vụ triển khai thực tế hiếm khi được khám phá. Khung giải pháp được đề xuất của chúng tôi lấp đầy khoảng trống này bằng cách thiết kế một Kiến trúc Ensemble lai mới (XGBoost, LightGBM, Random Forest) kết hợp với cân bằng dựa trên SMOTEENN, và sau đó mở rộng bằng kỹ thuật Chưng cất Tri thức (Knowledge Distillation). Phương pháp này chuyển đổi hiệu quả các ranh giới quyết định phức tạp của một mô hình ensemble đồ sộ sang một mô hình Student nhẹ, đảm bảo cả độ chính xác cao và độ trễ suy luận phù hợp với môi trường điện toán biên.
